\documentclass[12pt]{article}

\usepackage[utf8]{inputenc}
\usepackage[T2A]{fontenc}
\usepackage[russian]{babel}
\usepackage{graphicx}
\usepackage[top=1.5cm, bottom=2cm, left=2cm, right=2cm, includehead, a4paper]{geometry}
\usepackage{amsmath}
\usepackage{amssymb}
\usepackage{tikz}
\usepackage{tabularx}
\usepackage{array} % ОБЯЗАТЕЛЬНО для таблиц с параметром m{} и C{}
\usepackage{hyperref}
\usepackage{fancyhdr}
\usepackage{pgffor} 

\usetikzlibrary{quotes, angles, calc, decorations.markings}

% --- НАСТРОЙКИ ДЛЯ ВАШЕЙ ТАБЛИЦЫ БАЛЛОВ ---
\newlength{\GFrowwidth}
\setlength{\GFrowwidth}{0.8cm}
\newlength{\hhhh}
\setlength{\hhhh}{5pt}
\setlength{\arrayrulewidth}{1.5pt}
\newcolumntype{C}[1]{>{\centering\arraybackslash}m{#1}}
% ------------------------------------------

% --- ЕДИНЫЙ БЛОК НАСТРОЕК КОЛОНТИТУЛОВ ---
\pagestyle{fancy}
\fancyhf{} 
\renewcommand{\headrulewidth}{0pt}
\setlength{\headheight}{15pt}
\setlength{\headsep}{15pt}


\newcommand{\grid}[2]{%
\begin{tikzpicture}[line cap=round,line join=round,x=0.5cm,y=0.5cm]
    \draw [xstep=0.5cm,ystep=0.5cm, gray!60, line width=0.2pt] (0.,0.) grid (#1, #2);
    % \clip(0.,0.) rectangle (3.,3.);
\end{tikzpicture}%
}



\NewDocumentCommand\samerowans{O{0.65}O{0.3}m}{\noindent\begin{minipage}{#1\linewidth}
#3
\end{minipage}%
\hfill\fbox{\begin{minipage}{#2\linewidth}
{\small\tt Ответ:}\raisebox{-5ex}{\rule{0pt}{7ex}}
\end{minipage}}}

\newcommand{\rowans}[1]{\par\noindent{\fbox{\begin{minipage}{.75\linewidth}
{\small\tt Ответ:}\raisebox{-5ex}{\rule{0pt}{7ex}}
\rule{0pt}{#1}
\end{minipage}}}}

\newenvironment{problem}[1]{\medskip\noindent\textbf{#1.}}{}


\newcommand\samerowpic[4]{\noindent\begin{minipage}{#3\textwidth}
#1
\end{minipage}%
\begin{minipage}{#4\textwidth}
\includegraphics[width=.9\textwidth]{#2}
\end{minipage}}

\newcommand\mytitle{\begin{center}
\textsc{Школа №1514. Экзаменационная работа по математике \\ для поступающих в 5 класс}
\medskip
\texttt{Демонстрационный вариант 4, 01 марта 2026. Продолжительность 60 минут}
\end{center}}

\begin{document}

\begin{center}
    \textbf{\large Вариант 1}
\end{center}

\samerowans{\begin{problem}{1}
Вычислите: $5709 - 5632 : (258 - 13 \cdot 19) + 103$
\end{problem}}

\vspace{5pt}

\samerowans{\begin{problem}{2}
Решите уравнение: $(294 : a - 14 \cdot 6) : 3 = 21$
\end{problem}}

\vspace{5pt}

\samerowans{\begin{problem}{3}
Вычислите: $9\text{ т } 54\text{ кг } - 76\text{ ц } + 127\text{ кг}$. Ответ запишите в т, ц, кг.
\end{problem}}

\samerowans{\begin{problem}{4}
Сложили пять одинаковых чисел, к результату прибавили 8, затем полученную сумму разделили на 9 и получили 12. Какие числа складывали?
\end{problem}}

\vspace{5pt}

\samerowans{\begin{problem}{5}
Оля два часа делала домашнее задание. Три четверти часа она делала математику, полчаса русский. Все оставшееся время она учила английский. Сколько минут Оля учила английский?
\end{problem}}

\vspace{5pt}

\begin{problem}{6} Из двух пунктов А и В, расстояние между которыми 50 км, одновременно навстречу друг другу выехали два велосипедиста. Скорость велосипедиста, выехавшего из п А, на $5\text{ км/ч}$ больше скорости второго. За два часа им удалось встретиться и разъехаться на расстояние 40 км. На каком расстоянии от п А будет второй велосипедист через 3 часа 15 минут после своего выезда из В?
\end{problem}

\vspace{5pt}

\noindent\grid{36}{19}

\vspace{5pt}

\samerowans{\begin{problem}{7}
Первый станок может сделать 300 деталей за 6 часов, второй справится с этой задачей в 2 раза быстрее. Сколько часов совместной работы потребуется двум станкам, чтобы сделать 750 деталей?
\end{problem}}

\newpage

\begin{problem}{8} Из прямоугольника с периметром 16 дм вырезали квадрат со стороной 5 см. Найди площадь оставшейся фигуры, если длина прямоугольника 500 мм.
\end{problem}

\smallskip

\noindent\grid{36}{15}

\vspace{5pt}

\samerowans{\begin{problem}{9}
Таня заплатила за покупку 1 м 50 см шелковой ленты 600 рублей. Надя купила ленты на 80 см больше, чем Таня. Сколько рублей заплатила Надя?
\end{problem}}

\vspace{5pt}

\samerowans{\begin{problem}{10}
Коля и Петя сыграли партию в шахматы, которая закончилась победой одного из них. «Я победил», — сказал Коля. «Я проиграл», — сказал Петя. Кто из них выиграл, если известно, что хотя бы кто-то из них сказал неправду?
\end{problem}}

\vspace{5pt}

\begin{problem}{11} Через два года мальчик будет вдвое старше, чем он был два года назад. А девочка будет через три года втрое старше, чем три года назад. Кто старше: мальчик или девочка?
\end{problem}

\smallskip

\noindent\grid{36}{20}

\newpage


\begin{center}
\textsc{\large Дополнительный лист}
\end{center}
\noindent\grid{34}{47}

\newpage

\null


\newpage

\begin{center}
    \textbf{\large Вариант 2}
\end{center}

\newpage
\null
\newpage

\samerowans{\begin{problem}{1}
Вычислите: $7567 - 7245 : (416 - 15 \cdot 17) + 78$
\end{problem}}

\vspace{5pt}

\samerowans{\begin{problem}{2}
Решите уравнение: $(24 \cdot 4 - 252 : b) \cdot 4 = 48$
\end{problem}}

\vspace{5pt}

\samerowans{\begin{problem}{3}
Вычислите: $8\text{ т } 45\text{ кг } - 54\text{ ц } + 136\text{ кг}$. Ответ запишите в т, ц, кг.
\end{problem}}

\samerowans{\begin{problem}{4}
Сложили пять одинаковых чисел, к результату прибавили 3, затем полученную сумму разделили на 9 и получили 17. Какие числа складывали?
\end{problem}}

\vspace{5pt}

\samerowans{\begin{problem}{5}
Юля два часа делала домашнее задание. Две трети часа она делала математику, полчаса русский. Все оставшееся время она учила английский. Сколько минут Юля учила английский?
\end{problem}}

\vspace{5pt}

\begin{problem}{6} Из двух пунктов А и В, расстояние между которыми 40 км, одновременно навстречу друг другу выехали два велосипедиста. Скорость велосипедиста, выехавшего из п А, на $5\text{ км/ч}$ больше скорости второго. За два часа им удалось встретиться и разъехаться на расстояние 30 км. На каком расстоянии от п А будет второй велосипедист через 3 часа 20 минут после своего выезда из В?
\end{problem}

\vspace{5pt}

\noindent\grid{36}{19}

\vspace{5pt}

\samerowans{\begin{problem}{7}
Первый станок может сделать 320 деталей за 8 часов, второй справится с этой задачей в 2 раза быстрее. Сколько часов совместной работы потребуется двум станкам, чтобы сделать 720 деталей?
\end{problem}}

\newpage

\begin{problem}{8} Из прямоугольника с периметром 20 дм вырезали квадрат со стороной 6 см. Найди площадь оставшейся фигуры, если длина прямоугольника 600 мм.
\end{problem}

\smallskip

\noindent\grid{36}{15}

\vspace{5pt}

\samerowans{\begin{problem}{9}
Таня заплатила за покупку 1 м 50 см шелковой ленты 750 рублей. Надя купила ленты на 40 см больше, чем Таня. Сколько рублей заплатила Надя?
\end{problem}}

\vspace{5pt}

\samerowans{\begin{problem}{10}
Коля и Петя сыграли партию в шахматы, которая закончилась победой одного из них. «Петя победил», — сказал Коля. «Коля проиграл», — сказал Петя. Кто из них выиграл, если известно, что хотя бы кто-то из них сказал неправду?
\end{problem}}

\vspace{5pt}

\begin{problem}{11} Через четыре года мальчик будет вдвое старше, чем он был четыре года назад. А девочка будет через три года втрое старше, чем семь лет назад. Кто старше: мальчик или девочка?
\end{problem}

\smallskip

\noindent\grid{36}{19}

\newpage



\newpage

\begin{center}
    \textbf{\large Вариант 3}
\end{center}

\samerowans{\begin{problem}{1}
Вычислите: $5104-4354 : (252-14\cdot17)+108$
\end{problem}}

\vspace{5pt}

\samerowans{\begin{problem}{2}
Решите уравнение: $(360 : a - 2 \cdot 15) \cdot 5 = 150$
\end{problem}}

\vspace{5pt}

\samerowans{\begin{problem}{3}
Вычислите: $6\text{ т } 84\text{ кг } - 45\text{ ц } + 126\text{ кг}$. Ответ запишите в т, ц, кг.
\end{problem}}

\samerowans{\begin{problem}{4}
Сложили пять одинаковых чисел, к результату прибавили 10, затем полученную сумму разделили на 8 и получили 15. Какие числа складывали?
\end{problem}}

\vspace{5pt}

\samerowans{\begin{problem}{5}
Катя два часа делала домашнее задание. Две трети часа она делала математику, одну четверть часа — русский язык. Все оставшееся время она учила английский. Сколько минут Катя учила английский?
\end{problem}}

\vspace{5pt}

\begin{problem}{6} Из двух пунктов А и В, расстояние между которыми 60 км, одновременно навстречу друг другу выехали два велосипедиста. Скорость велосипедиста, выехавшего из п А, на $5\text{ км/ч}$ больше скорости второго. За два часа им удалось встретиться и разъехаться на расстояние 30 км. На каком расстоянии от п А будет второй велосипедист через 4 часа 15 минут после своего выезда из В?
\end{problem}

\vspace{5pt}

\noindent\grid{36}{19}

\vspace{5pt}

\samerowans{\begin{problem}{7}
Первый станок может сделать 420 деталей за 6 часов, второй справится с этой задачей в 2 раза быстрее. Сколько часов совместной работы потребуется двум станкам, чтобы сделать 630 деталей?
\end{problem}}

\newpage

\begin{problem}{8} Из прямоугольника с периметром 22 дм вырезали квадрат со стороной 4 см. Найди площадь оставшейся фигуры, если длина прямоугольника 700 мм.
\end{problem}

\smallskip

\noindent\grid{36}{15}

\vspace{5pt}

\samerowans{\begin{problem}{9}
Таня заплатила за покупку 1 м 40 см шелковой ленты 420 рублей. Надя купила ленты на 50 см больше, чем Таня. Сколько рублей заплатила Надя?
\end{problem}}

\vspace{5pt}

\samerowans{\begin{problem}{10}
Коля и Петя сыграли партию в шахматы, которая закончилась победой одного из них. «Петя выиграл», — сказал Коля. «Я победил», — сказал Петя. Кто из них выиграл, если известно, что хотя бы кто-то из них сказал неправду?
\end{problem}}

\vspace{5pt}

\begin{problem}{11} Через 8 лет мальчик будет вдвое старше, чем он был 3 года назад. А девочка будет через 4 года втрое старше, чем 8 лет назад. Кто старше: мальчик или девочка?
\end{problem}

\smallskip

\noindent\grid{36}{20}

\newpage

\begin{center}
\large{Вариант 3}
\end{center}

\begin{enumerate}
    \item \textbf{Вычислите:} $5104 - 4354 : (252 - 14 \cdot 17) + 108$ \\
    \textit{Решение:}
    \begin{itemize}
        \item $14 \cdot 17 = 238$
        \item $252 - 238 = 14$
        \item $4354 : 14 = 311$
        \item $5104 - 311 + 108 = 4901$
    \end{itemize}
    \textbf{Ответ: 4901}

    \item \textbf{Решите уравнение:} $(360:a - 2 \cdot 15) \cdot 5 = 150$ \\
    \textit{Решение:}
    \begin{itemize}
        \item $360:a - 30 = 150 : 5$
        \item $360:a - 30 = 30$
        \item $360:a = 60 \implies a = 6$
    \end{itemize}
    \textbf{Ответ: 6}

    \item \textbf{Вычислите:} 6 т 84 кг $- 45$ ц $+ 126$ кг \\
    \textit{Решение:}
    \begin{itemize}
        \item $6084 \text{ кг} - 4500 \text{ кг} + 126 \text{ кг} = 1710 \text{ кг}$
        \item $1710 \text{ кг} = 1 \text{ т } 7 \text{ ц } 10 \text{ кг}$
    \end{itemize}
    \textbf{Ответ: 1 т 7 ц 10 кг}

    \item \textbf{Задача (пять одинаковых чисел):} \\
    \textit{Решение:} Пусть $x$ — искомое число. \\
    $(5x + 10) : 8 = 15 \implies 5x + 10 = 120 \implies 5x = 110 \implies x = 22$ \\
    \textbf{Ответ: 22}

    \item \textbf{Домашнее задание Кати:} \\
    \textit{Решение:} Общее время 2 часа = 120 мин.
    \begin{itemize}
        \item Математика: $\frac{2}{3} \cdot 60 = 40$ мин
        \item Русский: $\frac{1}{4} \cdot 60 = 15$ мин
        \item Английский: $120 - (40 + 15) = 65$ мин
    \end{itemize}
    \textbf{Ответ: 65 минут}

    \item \textbf{Велосипедисты:} \\
    \textit{Решение:}
    \begin{itemize}
        \item Суммарный путь: $60 + 30 = 90$ км. Скорость сближения: $90 / 2 = 45$ км/ч.
        \item $v_1 + v_2 = 45$ и $v_1 = v_2 + 5 \implies v_2 = 20, v_1 = 25$ км/ч.
        \item Путь второго из B за 1 ч 15 мин: $20 \cdot 1.25 = 25$ км.
        \item Расстояние от A: $60 - 25 = 35$ км.
    \end{itemize}
    \textbf{Ответ: 35 км}

    \item \textbf{Станки:} \\
    \textit{Решение:}
    \begin{itemize}
        \item Производительность 1: $420 / 6 = 70$ дет/ч. Производительность 2: $420 / 3 = 140$ дет/ч.
        \item Время совместной работы: $630 / (70 + 140) = 3$ часа.
    \end{itemize}
    \textbf{Ответ: 3 часа}

    \item \textbf{Площадь фигуры:} \\
    \textit{Решение:} $P = 220$ см, длина $a = 70$ см.
    \begin{itemize}
        \item Ширина $b = (220 - 2 \cdot 70) / 2 = 40$ см.
        \item $S = (70 \cdot 40) - (4 \cdot 4) = 2800 - 16 = 2784 \text{ см}^2$.
    \end{itemize}
    \textbf{Ответ: 2784 см$^2$}

    \item \textbf{Шелковая лента:} \\
    \textit{Решение:} 140 см стоят 420 руб $\implies$ 1 см стоит 3 руб. \\
    Надя купила $140 + 50 = 190$ см. Стоимость: $190 \cdot 3 = 570$ руб. \\
    \textbf{Ответ: 570 рублей}

    \item \textbf{Шахматы (Логика):} \\
    \textit{Решение:} Оба утверждения означали победу Пети. Так как по условию один из них солгал, победа Пети невозможна. Выиграл Коля. \\
    \textbf{Ответ: Коля}

    \item \textbf{Возраст:} \\
    \textit{Решение:}
    \begin{itemize}
        \item Мальчик: $x + 8 = 2(x - 3) \implies x = 14$ лет.
        \item Девочка: $y + 4 = 3(y - 8) \implies y = 14$ лет.
    \end{itemize}
    \textbf{Ответ: одинаково (по 14 лет)}
\end{enumerate}

\newpage

\begin{center}
    \textbf{\large Вариант 4}
\end{center}

\samerowans{\begin{problem}{1}
Вычислите: $3816 - 3640 : (330 - 16 \cdot 19) + 84$
\end{problem}}

\vspace{5pt}

\samerowans{\begin{problem}{2}
Решите уравнение: $(18 \cdot 4 - 450 : b) : 3 = 18$
\end{problem}}

\vspace{5pt}

\samerowans{\begin{problem}{3}
Вычислите: $7\text{ т } 56\text{ кг } - 38\text{ ц } + 244\text{ кг}$. Ответ запишите в т, ц, кг.
\end{problem}}

\samerowans{\begin{problem}{4}
Сложили пять одинаковых чисел, из результата вычли 12, затем полученную разность разделили на 6 и получили 13. Какие числа складывали?
\end{problem}}

\vspace{5pt}

\samerowans{\begin{problem}{5}
Миша два часа делал домашнее задание. Три пятых часа он делал математику, одну шестую часа — русский язык. Все оставшееся время он учил английский. Сколько минут Миша учил английский?
\end{problem}}

\vspace{5pt}

\begin{problem}{6} Из двух пунктов А и В, расстояние между которыми 80 км, одновременно навстречу друг другу выехали два экскаватора. Скорость первого экскаватора, выехавшего из п. А, на $5\text{ км/ч}$ больше скорости второго. За два часа им удалось встретиться и разъехаться на расстояние 50 км. На каком расстоянии от п А будет второй экскаватор через 3 часа 20 минут после своего выезда из В?
\end{problem}

\vspace{5pt}

\noindent\grid{36}{19}

\vspace{5pt}

\samerowans{\begin{problem}{7}
Первый станок может сделать 480 деталей за 8 часов, второй справится с этой задачей в 2 раза быстрее. Сколько часов совместной работы потребуется двум станкам, чтобы сделать 540 деталей?
\end{problem}}

\newpage

\begin{problem}{8} Из прямоугольника с периметром 18 дм вырезали квадрат со стороной 3 см. Найди площадь оставшейся фигуры, если ширина прямоугольника 400 мм.
\end{problem}

\smallskip

\noindent\grid{36}{15}

\vspace{5pt}

\samerowans{\begin{problem}{9}
Таня заплатила за покупку 1 м 40 см шелковой ленты 560 рублей. Надя купила ленты на 50 см больше, чем Таня. Сколько рублей заплатила Надя?
\end{problem}}

\vspace{5pt}

\samerowans{\begin{problem}{10}
Коля и Петя сыграли партию в шахматы, которая закончилась победой одного из них. «Я победил», — сказал Коля. «Коля не проиграл», — сказал Петя. Кто из них выиграл, если известно, что хотя бы кто-то из них сказал неправду?
\end{problem}}

\vspace{5pt}

\begin{problem}{11} Через шесть лет мальчик будет втрое старше, чем он был два года назад. А девочка будет через два года вдвое старше, чем пять лет назад. Кто старше: мальчик или девочка?
\end{problem}

\smallskip

\noindent\grid{36}{20}

\newpage

\begin{center}
    \large{Вариант 4}
\end{center}


\begin{enumerate}
    \item \textbf{Вычислите:} $3816 - 3640 : (330 - 16 \cdot 19) + 84$ \\
    \textit{Решение:}
    \begin{itemize}
        \item $16 \cdot 19 = 304$
        \item $330 - 304 = 26$
        \item $3640 : 26 = 140$
        \item $3816 - 140 + 84 = 3760$
    \end{itemize}
    \textbf{Ответ: 3760}

    \item \textbf{Решите уравнение:} $(18 \cdot 4 - 450 : b) : 3 = 18$ \\
    \textit{Решение:}
    \begin{itemize}
        \item $72 - 450 : b = 54$
        \item $450 : b = 18 \implies b = 25$
    \end{itemize}
    \textbf{Ответ: 25}

    \item \textbf{Вычислите:} 7 т 56 кг $- 38$ ц $+ 244$ кг \\
    \textit{Решение:} $7056 \text{ кг} - 3800 \text{ кг} + 244 \text{ кг} = 3500 \text{ кг} = 3 \text{ т } 5 \text{ ц } 0 \text{ кг}$. \\
    \textbf{Ответ: 3 т 5 ц 0 кг}

    \item \textbf{Задача (пять одинаковых чисел):} \\
    \textit{Решение:} $(5x - 12) : 6 = 13 \implies 5x = 90 \implies x = 18$. \\
    \textbf{Ответ: 18}

    \item \textbf{Домашнее задание Миши:} \\
    \textit{Решение:} 2 часа = 120 мин.
    \begin{itemize}
        \item Математика: 36 мин, Русский: 10 мин.
        \item Английский: $120 - 46 = 74$ мин.
    \end{itemize}
    \textbf{Ответ: 74 минуты}

    \item \textbf{Экскаваторы:} \\
    \textit{Решение:}
    \begin{itemize}
        \item Суммарный путь: $80 + 50 = 130$ км. Скорость сближения: 65 км/ч.
        \item $v_1 = v_2 + 5 \implies v_2 = 30, v_1 = 35$ км/ч.
        \item Путь второго из B за $3\frac{1}{3}$ ч: $30 \cdot 10/3 = 100$ км.
        \item Расстояние от A: $|100 - 80| = 20$ км.
    \end{itemize}
    \textbf{Ответ: 20 км}

    \item \textbf{Станки:} \\
    \textit{Решение:}
    \begin{itemize}
        \item Пр1: 60 дет/ч, Пр2: 120 дет/ч.
        \item Время вместе: $540 / 180 = 3$ часа.
    \end{itemize}
    \textbf{Ответ: 3 часа}

    \item \textbf{Площадь фигуры:} \\
    \textit{Решение:} $P = 180$ см, ширина $b = 40$ см $\implies$ длина $a = 50$ см. \\
    $S = (50 \cdot 40) - (3 \cdot 3) = 1991 \text{ см}^2$. \\
    \textbf{Ответ: 1991 см$^2$ }

    \item \textbf{Шелковая лента:} \\
    \textit{Решение:} 1 см стоит 4 руб. Надя купила 190 см $\implies$ $190 \cdot 4 = 760$ руб. \\
    \textbf{Ответ: 760 рублей}

    \item \textbf{Шахматы (Логика):} \\
    \textit{Решение:} Оба утверждения означают победу Коли. Так как кто-то из них солгал, победил Петя. \\
    \textbf{Ответ: Петя}

    \item \textbf{Возраст:} \\
    \textit{Решение:}
    \begin{itemize}
        \item Мальчик: $x + 6 = 3(x - 2) \implies x = 6$ лет.
        \item Девочка: $y + 2 = 2(y - 5) \implies y = 12$ лет.
    \end{itemize}
    \textbf{Ответ: девочка (мальчику 6 лет, девочке 12 лет)}
\end{enumerate}

\end{document}